% © 2026 Ing. David Jaroš — CC BY-NC-ND 4.0
%
% This work is licensed under a Creative Commons Attribution-NonCommercial-NoDerivatives
% 4.0 International License (CC BY-NC-ND 4.0).
%
% License History: Earlier drafts (up to v0.3) were released under CC BY 4.0.
% From v0.4 onward, all material is released under CC BY-NC-ND 4.0 to protect
% the integrity of the theoretical work during ongoing academic development.
%
% See LICENSE.md for full license text.
%
% Gap ID: Gap C2 Step 1
% Purpose: Derive SM fermion hypercharge assignments Y from UBT ψ-winding
%          topology and the OP-S4 decoupling of SU(2)_R.
% Status: [MC/OPEN] — Y=(B-L)/2 formula established [L1 cond. on OP-S4];
%         sub-gap C2-i (fractional quark winding) remains [OPEN/MC].
% Date: 2026-06-11

\ifdefined\INCLUDEMODE
\else
\documentclass[12pt,a4paper]{article}
\usepackage[a4paper,margin=2.5cm]{geometry}
\usepackage{amsmath,amssymb,amsthm}
\usepackage{mathtools}
\usepackage{hyperref}
\usepackage{xcolor}
\usepackage{mdframed}
\usepackage{booktabs}

\hypersetup{
  colorlinks=true,
  linkcolor=blue!60!black,
  citecolor=green!50!black,
  urlcolor=blue!60!black
}

\newtheorem{theorem}{Theorem}[section]
\newtheorem{lemma}[theorem]{Lemma}
\newtheorem{proposition}[theorem]{Proposition}
\newtheorem{corollary}[theorem]{Corollary}
\theoremstyle{definition}
\newtheorem{definition}[theorem]{Definition}
\theoremstyle{remark}
\newtheorem{remark}[theorem]{Remark}

\newmdenv[backgroundcolor=yellow!20, linecolor=orange, linewidth=1.5pt]{gapbox}
\newmdenv[backgroundcolor=green!15, linecolor=green!60!black, linewidth=1.5pt]{resultbox}
\newmdenv[backgroundcolor=red!10, linecolor=red!60, linewidth=1.5pt]{deadendbox}
\newmdenv[backgroundcolor=blue!8, linecolor=blue!50, linewidth=1.5pt]{insightbox}

\title{\textbf{Gap C2 Step 1: Fermion Hypercharge Assignments from UBT}\\
       \large $Y = (B{-}L)/2$ from $\psi$-Winding Topology and $SU(2)_R$ Decoupling}
\author{Ing. David Jaroš}
\date{June 2026}

\begin{document}
\maketitle

\begin{abstract}
We derive the Standard Model hypercharge assignments
$Y(Q_L)=\tfrac{1}{6}$, $Y(u_R)=\tfrac{2}{3}$, $Y(d_R)=-\tfrac{1}{3}$,
$Y(L)=-\tfrac{1}{2}$, $Y(e_R)=-1$, $Y(\nu_R)=0$
from UBT $\psi$-winding topology under the assumption that $SU(2)_R$ is
decoupled (OP-S4 [L1 conditional]).
The key relation $Y = (B{-}L)/2$ is established at
\textbf{[L1 conditional on OP-S4 + sub-gap C2-i]}.
A remaining sub-gap (C2-i) concerns the first-principles derivation of
fractional baryon-winding numbers for quarks from the UBT SU(3) action.
\end{abstract}

\fi

%─────────────────────────────────────────────────────────────────────────────
\section{Background and Chain Position}
\label{sec:background}
%─────────────────────────────────────────────────────────────────────────────

This file records Gap~C2 Step~1.
The T2\_GAUGE paper states $\sin^2\theta_W$ at
\textbf{[L1 conditional on Gap~C2]}.
Closing Gap~C2 Step~1 requires deriving the SM hypercharge assignments
for all six fermion representations from UBT without fitting them to experiment.

\textbf{Input from prior steps:}
\begin{itemize}
  \item $U(1)_Y$ arises as the right scalar phase action
        $\Theta \mapsto e^{-i\theta}\Theta$
        (established [L0] in \texttt{canonical/interactions/sm\_gauge.tex}).
        This identifies $U(1)_Y$ as a symmetry \emph{group} but does not
        yet assign specific $Y$ values to fermion representations.
  \item $SU(2)_R$ is decoupled from the low-energy spectrum (OP-S4 [L1 conditional])
        via $\psi$-parity and Kaluza-Klein mass suppression.
        Consequently, the right-handed weak isospin $T_{3R}=0$ for all
        observed fermions.
  \item Fermion generations are labelled by the $\psi$-winding number $n$
        (three generations from $\psi$-winding [L0] in
        \texttt{canonical/interactions/sm\_gauge.tex}).
\end{itemize}

%─────────────────────────────────────────────────────────────────────────────
\section{The Hypercharge Formula $Y = (B-L)/2$}
\label{sec:hypercharge_formula}
%─────────────────────────────────────────────────────────────────────────────

\begin{lemma}[Hypercharge from $\psi$-winding and OP-S4,
              \textbf{[L1 conditional on OP-S4 + sub-gap C2-i]}]
\label{lem:hypercharge_formula}
Assume:
\begin{enumerate}
  \item[(A)] $SU(2)_R$ is decoupled at the observational energy scale
             (OP-S4 [L1 conditional]), so the right-handed weak isospin
             satisfies $T_{3R}=0$ for every observed fermion.
  \item[(B)] The $U(1)_Y$ charge is embedded in the extended electroweak
             group as
             \begin{equation}
               Y \;=\; T_{3R} + \frac{B{-}L}{2},
               \label{eq:gell_mann_nishijima}
             \end{equation}
             where $B$ is baryon number and $L$ is lepton number.
  \item[(C)] Baryon and lepton numbers are assigned by $\psi$-winding topology:
             quarks carry $B=1/3$ per quark (hence $B=1$ for a proton) and
             leptons carry $L=1$.
\end{enumerate}
Then, under assumption (A), equation \eqref{eq:gell_mann_nishijima} reduces to
\begin{equation}
  \boxed{Y \;=\; \frac{B{-}L}{2},}
  \label{eq:Y_BL}
\end{equation}
and all six SM hypercharge assignments follow algebraically.
\end{lemma}

\begin{proof}
Set $T_{3R}=0$ in \eqref{eq:gell_mann_nishijima}.
The baryon-minus-lepton charge $B{-}L$ and resulting hypercharge
$Y = T_{3R} + (B{-}L)/2$ for each representation are:
\begin{center}
\begin{tabular}{lccccc}
\toprule
Representation & $B$ & $L$ & $B{-}L$ & $T_{3R}$ (pre-OP-S4) & $Y$ \\
\midrule
$Q_L = (u_L, d_L)^T$ & $1/3$ & $0$ & $1/3$ & $0$ & $1/6$ \\
$u_R$               & $1/3$ & $0$ & $1/3$ & $+1/2$ & $2/3$ \\
$d_R$               & $1/3$ & $0$ & $1/3$ & $-1/2$ & $-1/3$ \\
$L   = (\nu_L, e_L)^T$ & $0$ & $1$ & $-1$ & $0$ & $-1/2$ \\
$e_R$               & $0$ & $1$ & $-1$ & $-1/2$ & $-1$ \\
$\nu_R$             & $0$ & $1$ & $-1$ & $+1/2$ & $0$ \\
\bottomrule
\end{tabular}
\end{center}

\noindent
The $T_{3R}$ column gives the right-handed weak isospin values in the full
$SU(2)_R$ theory; these are used to compute the $Y$ values via
$Y = T_{3R} + (B{-}L)/2$.
The hypercharge values so obtained are group-theoretic properties of the
representations, fixed before the OP-S4 decoupling.
Applying OP-S4 (assumption~(A)) decouples the \emph{gauge coupling} to $SU(2)_R$
but does not change the $Y$ eigenvalues, since $Y$ is a generator of the
commuting $U(1)_Y$ factor.
All six values match the Standard Model hypercharge table.
\end{proof}

\begin{remark}[Consistency check via Gell-Mann--Nishijima]
\label{rem:gm_nishijima}
The electric charge formula $Q = T_{3L} + Y$ reproduces the SM charges:
$Q(u_L) = 1/2 + 1/6 = 2/3$, $Q(d_L) = -1/2 + 1/6 = -1/3$,
$Q(e_L) = -1/2 - 1/2 = -1$, $Q(\nu_L) = 1/2 - 1/2 = 0$.
This is a consistency check, not a new derivation.
\end{remark}

%─────────────────────────────────────────────────────────────────────────────
\section{Residual Sub-Gap C2-i: Fractional Quark Winding}
\label{sec:c2i}
%─────────────────────────────────────────────────────────────────────────────

Assumption (C) in Lemma~\ref{lem:hypercharge_formula} asserts $B=1/3$ per quark
from $\psi$-winding topology.
In the UBT SU(3) derivation (\texttt{canonical/su3\_derivation/}), quarks sit in
the $\mathbf{3}$ representation of $SU(3)_c$, which is associated with
$\mathbb{Z}_2{\times}\mathbb{Z}_2{\times}\mathbb{Z}_2$ involutions on $\mathbb{C}\otimes\mathbb{H}$.
However, the fractional winding number $n_q = 1/3$ for quarks (versus $n_\ell = 1$
for leptons) has not been derived from first principles in UBT.

\begin{gapbox}
  \textbf{Sub-gap C2-i [OPEN/MC]:} Derive the assignment $B=1/3$ per quark
  (equivalently, $\psi$-winding number $n_q = 1/3$) from the SU(3) colour-charge
  structure in UBT without using it as an input.

  \textbf{Candidate approach}: The colour-singlet constraint in SU(3) requires
  three quarks to combine into a baryon with $B=1$, suggesting $B_q = 1/3$ from
  the combinatorial structure of the $\mathbf{3}$-representation.
  A precise UBT statement would require showing that the $\psi$-winding eigenvalues
  of states in $\mathbf{3}$ of SU(3) are $1/3$ of those in $\mathbf{1}$ (singlet/lepton),
  following from the $\mathbb{Z}_2{\times}\mathbb{Z}_2{\times}\mathbb{Z}_2$ involution
  structure.

  \textbf{Status}: [OPEN/MC] — no current complete proof.
  Does not block T2\_GAUGE paper (hypercharge formula works conditionally on this sub-gap).
\end{gapbox}

%─────────────────────────────────────────────────────────────────────────────
\section{Status Summary}
\label{sec:status}
%─────────────────────────────────────────────────────────────────────────────

\begin{resultbox}
  \textbf{Gap C2 Step 1 outcome (2026-06-11):}

  \begin{center}
  \begin{tabular}{lll}
    \toprule
    \textbf{Claim} & \textbf{Level} & \textbf{Source} \\
    \midrule
    $Y = (B{-}L)/2$ from OP-S4 + $\psi$-winding & [L1 cond. on OP-S4 + C2-i] & Lem.~\ref{lem:hypercharge_formula} \\
    All 6 SM $Y$-values reproduced & [L0] (algebra given $Y=(B-L)/2$ and $T_{3R}$) & proof above \\
    Sub-gap C2-i: $B_q=1/3$ from SU(3) & [OPEN/MC] & §\ref{sec:c2i} \\
    \bottomrule
  \end{tabular}
  \end{center}

  The hypercharge formula $Y=(B{-}L)/2$ is obtained from the well-established
  embedding of $U(1)_Y \subset SU(5)$ (or equivalently, from $T_{3R}=0$ under OP-S4)
  combined with the $\psi$-winding B-L assignment.
  The remaining sub-gap C2-i is the first-principles derivation of $B_q=1/3$
  from the UBT SU(3) colour structure.

  \textbf{Impact on T2\_GAUGE}: Gap~C2 Step~1 is now partially closed.
  The anomaly-cancellation statement in the T2\_GAUGE paper §anomaly
  (referenced as [L1 conditional on C2]) can be updated to
  [L1 conditional on C2-i] — a strictly weaker and more precisely
  stated condition.
\end{resultbox}

\ifdefined\INCLUDEMODE
\else
\end{document}
\fi
